\documentclass[11pt]{article}
\usepackage[margin=1in]{geometry}
\usepackage{amsmath,amssymb,amsthm}
\usepackage{booktabs}
\usepackage{array}
\usepackage{xcolor}
\usepackage{hyperref}
\usepackage{parskip}
\usepackage{enumitem}
\usepackage{listings}
\usepackage{caption}

\hypersetup{colorlinks=true, linkcolor=blue, urlcolor=blue, citecolor=blue}

\lstset{
  basicstyle=\ttfamily\small,
  breaklines=true,
  frame=single,
  backgroundcolor=\color{gray!8},
  xleftmargin=6pt, xrightmargin=6pt
}

\title{\textbf{The \texttt{avg\_time\_periods} Option in \texttt{did\_multiplegt\_dyn}}\\[6pt]
\large Requirements, Calculations, and Per-Option Behaviour}
\author{Development notes --- wagepan cross-check}
\date{June 2026}

\begin{document}
\maketitle
\tableofcontents
\newpage

%───────────────────────────────────────────────────────────────────────────────
\section{Purpose}
%───────────────────────────────────────────────────────────────────────────────

The \texttt{avg\_time\_periods} option computes two quantities and exposes them
as \texttt{e()} scalars after a call to \texttt{did\_multiplegt\_dyn}:

\begin{enumerate}[nosep]
  \item \textbf{\texttt{e(avg\_cumul)}}: the average number of time periods over
        which a treatment effect is accumulated, weighted by the magnitude of
        each switcher's treatment change.
  \item \textbf{\texttt{e(N\_switch\_avg\_k)}} for $k = 1,\dots,L$: the number
        of switcher groups that contribute to the $k$-th dynamic effect, as
        seen by the \texttt{avg\_time\_periods} computation. These are required
        to equal \texttt{e(N\_switchers\_effect\_k)} from the standard pipeline
        for the option to be internally consistent.
\end{enumerate}

The computations are carried out by two Mata functions,
\texttt{\_compute\_Tg} and \texttt{\_avg\_cumul}, which are defined at the top
of \texttt{did\_multiplegt\_dyn\_reset.do} (Program 1, lines~57--307) and
called from the \texttt{avg\_time\_periods} Stata block (lines~3272--3326).

%───────────────────────────────────────────────────────────────────────────────
\section{Option Description (help-file entry)}
%───────────────────────────────────────────────────────────────────────────────

The following text is the description added to the \textbf{Other options}
section of \texttt{did\_multiplegt\_dyn\_reset.sthlp}.

\bigskip
\noindent\rule{\linewidth}{0.4pt}

\noindent\textbf{avg\_time\_periods}: when this option is specified, the command
reports the average number of time periods over which a switcher group's
treatment effect is accumulated, and the number of switchers entering the
estimation of each dynamic effect.

For each ever-switcher group $g$, let $F_g$ denote its date of first treatment
change, let $T_g$ denote the last period for which a valid control group exists
at every horizon from $F_g$ to $T_g$ (given the sample and the options
requested), and let $M_g = \min(L,\, T_g - F_g + 1)$, where $L$ is the number
of effects requested with \texttt{effects()}. The command returns the scalar
\texttt{e(avg\_cumul)}, defined as

\[
  \texttt{e(avg\_cumul)}
  \;=\;
  \frac{\displaystyle\sum_{g}\;\sum_{k=0}^{M_g-1}
        |D_{g,F_g+k} - D_{g,1}|\cdot(M_g - k)}
       {\displaystyle\sum_{g}\;\sum_{k=0}^{M_g-1}
        |D_{g,F_g+k} - D_{g,1}|},
\]

where the outer sums run over all switcher groups $g$ for which a valid control
group exists at horizon $k$, and $D_{g,1}$ is group $g$'s period-one
treatment. For a binary treatment, \texttt{e(avg\_cumul)} simplifies to a
weighted mean of the $M_g$ values with weights proportional to $M_g$. The
command also returns, for each $k = 1,\dots,L$, the scalar
\texttt{e(N\_switch\_avg\_k)}: the (weighted) number of switcher groups
contributing to dynamic effect $k$, which equals
\texttt{e(N\_switchers\_effect\_k)} from the standard pipeline.

\noindent\rule{\linewidth}{0.4pt}
\bigskip

%───────────────────────────────────────────────────────────────────────────────
\section{Notation}
%───────────────────────────────────────────────────────────────────────────────

\begin{center}
\renewcommand{\arraystretch}{1.3}
\begin{tabular}{@{}lp{10cm}@{}}
\toprule
Symbol & Meaning \\
\midrule
$g$ & Group identifier (\texttt{group\_XX}) \\
$t$ & Time period (\texttt{time\_XX}) \\
$D_{g,t}$ & Treatment of group $g$ at time $t$ (\texttt{treatment\_XX}) \\
$Y_{g,t}$ & Outcome of group $g$ at time $t$ (\texttt{outcome\_XX}) \\
$D_{g,1}$ & Baseline (period-1) treatment of group $g$ (\texttt{d\_sq\_XX};
             \texttt{d\_sq\_XX\_orig} when \texttt{continuous()}) \\
$F_g$ & First period when $D_{g,t} \neq D_{g,1}$; set to $T_{\max}+1$ for
        never-switchers (\texttt{F\_g\_XX}) \\
$T_{\max}$ & Last time period after trimming (\texttt{T\_max\_XX}) \\
$T_g^{\text{cls}}$ & Class-level ceiling: last period for which a control group
                     with the same baseline treatment (and same
                     \texttt{trends\_nonparam} level, if specified) exists.
                     $T_g^{\text{cls}} = \max_{g' \in \text{cls}(g)} F_{g'}^{\text{trunc}} - 1$,
                     stored in \texttt{T\_g\_XX} \\
$F_g^{\text{trunc}}$ & $\min(F_g,\, t_g^{\text{trunc}})$ where $t_g^{\text{trunc}}$ is
                        the last period at which at least one control exists for
                        group $g$'s class (\texttt{F\_g\_trunc\_XX}) \\
$T_g^{(h)}$ & Per-horizon last valid period for switcher $g$: the last $P \geq F_g$
              such that (i) both $g$ and a control group have outcome observed
              at $F_g-1$ and at $P$, and (ii) the control has not yet switched
              by $P$. Computed by \texttt{\_compute\_Tg}; stored in
              \texttt{Tg\_ph\_XX} \\
$M_g$ & Number of valid horizons for switcher $g$:
        $M_g = \min\!\bigl(L,\; T_g^{(h)} - F_g + 1\bigr)$,
        stored in \texttt{Mg\_ph\_XX} \\
$L$ & Number of dynamic effects requested (\texttt{effects()}, \texttt{l\_XX}) \\
$S_g$ & Switcher direction: $S_g = 1$ for upward switches, $S_g = 0$ for
        downward switches (\texttt{S\_g\_XX}) \\
$W_g$ & Cell weight of group $g$ at the relevant period (\texttt{N\_gt\_XX}) \\
$\text{cls}(g)$ & The strata of group $g$: groups sharing the same
                 \texttt{d\_sq\_int\_XX} value (integer code for $D_{g,1}$,
                 further stratified by \texttt{trends\_nonparam} when
                 \texttt{T\_g\_XX} is defined) \\
\bottomrule
\end{tabular}
\end{center}

%───────────────────────────────────────────────────────────────────────────────
\section{Data Preparation Before the \texttt{avg\_time\_periods} Block}
%───────────────────────────────────────────────────────────────────────────────

The \texttt{avg\_time\_periods} block (lines~3272--3326) runs \emph{after} the
main data-cleaning pipeline of Program~1 and \emph{after} the standard
estimator is computed by Program~2 (\texttt{did\_multiplegt\_dyn\_core\_new}).
By the time \texttt{avg\_time\_periods} executes, the working dataset already
contains the following key variables.

\subsection{Class Ceiling \texorpdfstring{$T_g^{\text{cls}}$}{Tg cls}
            (\texttt{T\_g\_XX})}

\begin{lstlisting}[language={}]
gen    F_g_trunc_XX = F_g_XX
replace F_g_trunc_XX = min(F_g_XX, trunc_control_XX) if trunc_control_XX != .
bys d_sq_XX `trends_nonparam': gegen T_g_XX = max(F_g_trunc_XX)
replace T_g_XX = T_g_XX - 1                    // line 1085-1086
\end{lstlisting}

$T_g^{\text{cls}}$ is the last period within a group's strata (same baseline
treatment, optionally same \texttt{trends\_nonparam} level) for which at least
one control group has not yet switched.

\subsection{Treatment Binarization for \texttt{continuous()} (line~1140--1143)}

When \texttt{continuous(p)} is specified, the continuous treatment is
binarized before the standard estimator runs:

\begin{lstlisting}[language={}]
gen treatment_XX_temp = (F_g_XX <= time_XX) * S_g_het_XX if S_g_het_XX != .
replace treatment_XX  = treatment_XX_temp
\end{lstlisting}

where $S_g^{\text{het}} = +1$ for upward switchers and $-1$ for downward
switchers. Crucially, \emph{never-switchers} have $S_g^{\text{het}} = \texttt{.}$,
so after this replacement \texttt{treatment\_XX $= $ .} for every observation
of a never-switcher. This has a direct implication for how the
\texttt{avg\_time\_periods} cell-validity indicator must be defined
(see Section~\ref{sec:ngtbin}).

\subsection{Baseline Treatment Variable (\texttt{d\_sq\_XX})}

For \texttt{continuous(p)}, the original continuous baseline is saved to
\texttt{d\_sq\_XX\_orig} before being replaced by zero:

\begin{lstlisting}[language={}]
gen  d_sq_XX_orig = d_sq_XX        // line 922
replace d_sq_XX = 0                // line 923
\end{lstlisting}

The \texttt{avg\_time\_periods} block uses \texttt{d\_sq\_XX\_orig} as the
baseline treatment in the weighted-average formula when \texttt{continuous()} is
active (\texttt{d1\_avg\_XX}; line~3277).

%───────────────────────────────────────────────────────────────────────────────
\section{Auxiliary Variables Created in the \texttt{avg\_time\_periods} Block}
%───────────────────────────────────────────────────────────────────────────────

Three temporary variables are created before the Mata calls.

\subsection{Ever-Switcher Indicator (\texttt{ever\_sw\_XX})}

\begin{lstlisting}[language={}]
gen byte `ever_sw_XX' = (F_g_XX < T_max_XX + 1)       // line 3282
\end{lstlisting}

A group is an ever-switcher if and only if its first switch period $F_g$ is
strictly before the sentinel value $T_{\max}+1$.

\subsection{Cell-Validity Indicator (\texttt{ngt\_bin\_XX})}
\label{sec:ngtbin}

\begin{lstlisting}[language={}]
gen byte `ngt_bin_XX' = (outcome_XX != .)              // line 3286
\end{lstlisting}

This is an unweighted binary flag that equals 1 whenever the outcome is
observed for that $(g,t)$ cell. It is the key indicator passed to the Mata
functions as the \texttt{NGT} column.

\textbf{Critical design choice.} The standard pipeline identifies control
$(g,t)$ cells via \texttt{diff\_y\_i\_XX $\neq$ .}, which requires only the
\emph{outcome} to be observed at two consecutive periods --- it does
\emph{not} require the treatment to be observed. When \texttt{continuous()} is
active, \texttt{treatment\_XX} is set to \texttt{.} for never-switchers after
binarization, but their outcomes remain available. Hence
\texttt{ngt\_bin\_XX} must be defined as \texttt{(outcome\_XX $\neq$ .)}
rather than \texttt{(outcome\_XX $\neq$ .\ \&\ treatment\_XX $\neq$ .)}.
The latter would mistakenly exclude never-switchers as controls in the
\texttt{continuous()} case, producing a systematic one-period offset in
\texttt{e(N\_switch\_avg\_k)} relative to \texttt{e(N\_switchers\_effect\_k)}.

\subsection{Per-Horizon Ceiling (\texttt{Tg\_ph\_XX})}

Initialized to \texttt{.} and then filled by \texttt{\_compute\_Tg}:

\begin{lstlisting}[language={}]
gen double `Tg_ph_XX' = .
sort group_XX time_XX
mata: _compute_Tg("group_XX","time_XX","outcome_XX","F_g_XX","T_g_XX",
      "`ever_sw_XX'","F_g_XX","d_sq_int_XX","`ngt_bin_XX'","`Tg_ph_XX'")
\end{lstlisting}

\subsection{Per-Horizon Horizon Count (\texttt{Mg\_ph\_XX})}

\begin{lstlisting}[language={}]
gen double `Mg_ph_XX' = min(`=l_XX', `Tg_ph_XX' - F_g_XX + 1)    // line 3301
\end{lstlisting}

$M_g$ is the number of post-switch periods over which switcher $g$ contributes,
capped at the requested number of effects $L$.

%───────────────────────────────────────────────────────────────────────────────
\section{The \texttt{\_compute\_Tg} Mata Function}
%───────────────────────────────────────────────────────────────────────────────

\subsection{Purpose}

For each ever-switcher $g$, \texttt{\_compute\_Tg} finds the per-horizon
ceiling $T_g^{(h)}$: the \emph{last} post-switch period $P$ for which (i)
group $g$ itself has a valid outcome at both the reference period $t^{\text{ref}}
= F_g - 1$ and at $P$, and (ii) at least one control group $g'$ (in the same
strata and with $F_{g'} > P$) also has valid outcomes at both periods.

\subsection{Algorithm (pseudo-code)}

\begin{lstlisting}[mathescape=true, language={}]
for each group g:
    Fg  = F_g_XX[g]
    Tgc = T_g_XX[g]       // class ceiling
    ever = ever_sw_XX[g]
    cls  = d_sq_int_XX[g]
    refp = Fg - 1          // FIXED reference period
    last_P = Fg - 1        // default: no valid horizon found

    if (ever == 1 and Fg < .) {
        own_ref = 1 iff NGT[g, refp] == 1 and Y[g, refp] < .

        if (own_ref == 1) {
            for P = Fg to Tgc {
                own_P = 1 iff NGT[g, P] == 1 and Y[g, P] < .

                if (own_P == 1) {
                    has_ctrl = 1 iff $\exists\, g' \neq g$: cls($g'$) == cls
                                        and $F_{g'} > P$
                                        and NGT[$g'$, refp] == 1 and Y[$g'$, refp] < .
                                        and NGT[$g'$, P]    == 1 and Y[$g'$, P]    < .
                    if (has_ctrl == 1): last_P = P
                }
                P++
            }
        }
    }
    Tg_ph_XX[g] = last_P   // written to all rows of group g
\end{lstlisting}

\subsection{Key Design Decisions}

\begin{itemize}[nosep]
  \item \textbf{Fixed reference period}: $t^{\text{ref}} = F_g - 1$ is
        computed once per group, outside the $P$-loop. This matches the
        standard pipeline, which uses $\Delta^L Y_{g,t} = Y_{g,t} -
        Y_{g,t-L}$ with $L = i$ evaluated at $t = F_g - 1 + i$, making
        the baseline always $F_g - 1$.

  \item \textbf{Control eligibility}: a potential control $g'$ must have
        $F_{g'} > P$ (has not yet switched by period $P$), must be in the
        same strata as $g$ (\texttt{d\_sq\_int\_XX}), and must have observed
        outcomes at both $t^{\text{ref}}$ and $P$.

  \item \textbf{Never-switchers as controls}: Since \texttt{NGT} is
        \texttt{ngt\_bin\_XX = (outcome != .)}, never-switchers (which have
        $F_{g'} = T_{\max}+1 > P$ for any feasible $P$) qualify as controls
        whenever their outcomes are observed --- exactly as in the standard
        pipeline.
\end{itemize}

%───────────────────────────────────────────────────────────────────────────────
\section{The \texttt{\_avg\_cumul} Mata Function}
%───────────────────────────────────────────────────────────────────────────────

\subsection{Purpose}

For each eligible switcher $g$ and each horizon $k = 0,\dots,M_g-1$,
\texttt{\_avg\_cumul} checks whether a valid control exists at the
\emph{current horizon} $\tau = F_g + k$, using the same fixed reference
$t^{\text{ref}} = F_g - 1$. If a control exists, the observation contributes
to the weighted average and to the per-effect switcher count.

\subsection{Eligibility Conditions}

A switcher group $g$ enters the outer loop only if all of the following hold:

\begin{align}
  & \texttt{ever} = 1 \text{ and } F_g < \infty \tag{is an ever-switcher}\\
  & F_g \leq T_g^{(h)} \tag{at least one valid horizon} \\
  & \texttt{ssw\_flag} = 0 \;\text{ or }\; M_g \geq L \tag{\texttt{same\_switchers}} \\
  & \texttt{sw\_dir} \neq \texttt{"in"} \;\text{ or }\; S_g = 1 \tag{\texttt{switchers(in)}} \\
  & \texttt{sw\_dir} \neq \texttt{"out"} \;\text{ or }\; S_g = 0 \tag{\texttt{switchers(out)}}
\end{align}

\subsection{Algorithm (pseudo-code)}

\begin{lstlisting}[mathescape=true, language={}]
num = 0;  den = 0;  nswitch = zeros(1, L)

for each eligible switcher g:
    refp       = Fg - 1                    // FIXED reference period
    y_ref      = Y[g, refp] if NGT[g,refp]==1, else missing
    d_ref_real = 1 iff y_ref is non-missing

    for k = 0 to Mg - 1:
        period = Fg + k

        // horizon observation check
        dval  = D[g, period] if NGT[g,period]==1 and Y[g,period] < .
        hor_ok = 1 iff dval, y_hor are non-missing

        if (dval < . and d_ref_real==1 and hor_ok==1 and y_ref < . and y_hor < .):
            has_ctrl = 1 iff $\exists\, g' \neq g$: cls($g'$)==cls
                                 and $F_{g'} > $ period
                                 and NGT[$g'$, refp]   ==1 and Y[$g'$, refp]   < .
                                 and NGT[$g'$, period] ==1 and Y[$g'$, period] < .

            if (has_ctrl == 1):
                incr   = |dval - d1|           // treatment magnitude
                weight = Mg - k                // periods still to accumulate
                num   += incr * weight
                den   += incr
                nswitch[k+1] += W[g]           // weighted switcher count
\end{lstlisting}

\subsection{Output Scalars}

\begin{align}
  \widehat{\overline{M}}
    &= \frac{\displaystyle\sum_{g \in \mathcal{S}}\;\sum_{k=0}^{M_g-1}
             \mathbf{1}[\text{ctrl}(g,k)]\cdot|D_{g,F_g+k} - D_{g,1}|\cdot(M_g - k)}
            {\displaystyle\sum_{g \in \mathcal{S}}\;\sum_{k=0}^{M_g-1}
             \mathbf{1}[\text{ctrl}(g,k)]\cdot|D_{g,F_g+k} - D_{g,1}|}
  \label{eq:avgM}
\end{align}

where $\mathcal{S}$ is the set of eligible switchers and $\mathbf{1}[\text{ctrl}(g,k)]$
is the indicator that a valid control group exists for horizon $k$ of switcher $g$.

The scalar $\widehat{\overline{M}}$ is stored in \texttt{avg\_cumul\_XX} and
subsequently posted as \texttt{e(avg\_cumul)}.

The row vector $\texttt{nswitch}[k]$ gives the weighted number of switchers
contributing to dynamic effect $k$; it is stored as Stata matrix
\texttt{nswitch\_mata\_XX} and later posted as \texttt{e(N\_switch\_avg\_k)}
for $k = 1,\dots,L$.

\subsection{Interpretation of the Weighted Average}

The weight $(M_g - k)$ assigned to each $(g, k)$ observation reflects the
number of remaining periods over which the treatment change of group $g$ is
still being accumulated when the $k$-th effect is estimated. For binary
treatment with $|D_{g,\tau} - D_{g,1}| = 1$ for all post-switch observations:

\[
  \widehat{\overline{M}}
  = \frac{\sum_g \sum_{k=0}^{M_g-1} (M_g - k)}{\sum_g M_g}
  = \frac{\sum_g \tfrac{M_g(M_g+1)}{2}}{\sum_g M_g}.
\]

This is a convex combination of the $M_g$ values, where groups with more valid
horizons receive relatively more weight.

%───────────────────────────────────────────────────────────────────────────────
\section{Final ereturn Scalars}
%───────────────────────────────────────────────────────────────────────────────

After the Mata calls, the per-effect counts are persisted as Stata scalars to
survive the \texttt{ereturn clear / ereturn post} that occurs later:

\begin{lstlisting}[language={}]
forvalues k = 1/`=l_XX' {
    scalar N_switch_avg_`k'_XX = nswitch_mata_XX[1, `k']
}
\end{lstlisting}

Further downstream, the following \texttt{e()} quantities are returned:

\begin{center}
\begin{tabular}{@{}ll@{}}
\toprule
\texttt{e()} scalar & Content \\
\midrule
\texttt{e(avg\_cumul)} & $\widehat{\overline{M}}$ from equation~\eqref{eq:avgM} \\
\texttt{e(N\_switch\_avg\_k)} & Weighted switcher count for effect $k$ \\
\bottomrule
\end{tabular}
\end{center}

The cross-check invariant is: for every valid specification,
$\texttt{e(N\_switch\_avg\_k)} = \texttt{e(N\_switchers\_effect\_k)}$.

%───────────────────────────────────────────────────────────────────────────────
\section{Per-Option Behaviour on the Wagepan Dataset}
%───────────────────────────────────────────────────────────────────────────────

The wagepan balanced panel (\texttt{nr}, \texttt{year}, \texttt{lwage},
\texttt{union}) is used to cross-check 19 option combinations with
\texttt{effects(5)} and, where applicable, \texttt{placebo(2)}.
All 19 specifications pass the invariant
$\texttt{e(N\_switch\_avg\_k)} = \texttt{e(N\_switchers\_effect\_k)}$
for every $k$.
Table~\ref{tab:specs} lists the specifications; the sections below describe how
each option interacts with the \texttt{avg\_time\_periods} computation.

\begin{table}[ht]
\centering
\caption{Wagepan specifications in \texttt{test\_wagepan\_specs.do}}
\label{tab:specs}
\renewcommand{\arraystretch}{1.2}
\begin{tabular}{@{}lp{9cm}@{}}
\toprule
Label & Options (always appended with \texttt{avg\_time\_periods graph\_off}) \\
\midrule
Baseline               & \texttt{effects(5)} \\
Placebos               & \texttt{effects(5) placebo(2)} \\
Normalized             & \texttt{effects(5) placebo(2) normalized} \\
Controls               & \texttt{effects(5) placebo(2) controls(hours)} \\
Trends\_Nonparam       & \texttt{effects(5) placebo(2) trends\_nonparam(black)} \\
Weight                 & \texttt{effects(5) placebo(2) weight(educ)} \\
Cluster                & \texttt{effects(5) placebo(2) cluster(hisp)} \\
Same\_Switchers        & \texttt{effects(5) placebo(2) same\_switchers} \\
Same\_Switchers\_Pl    & \texttt{effects(5) placebo(2) same\_switchers same\_switchers\_pl} \\
Switchers\_In          & \texttt{effects(5) placebo(2) switchers(in)} \\
Switchers\_Out         & \texttt{effects(5) placebo(2) switchers(out)} \\
Only\_Never\_Switchers & \texttt{effects(5) placebo(2) only\_never\_switchers} \\
CI\_Level\_90          & \texttt{effects(5) placebo(2) ci\_level(90)} \\
CI\_Level\_99          & \texttt{effects(5) placebo(2) ci\_level(99)} \\
Less\_Conservative\_SE & \texttt{effects(5) placebo(2) less\_conservative\_se} \\
Dont\_Drop\_Larger\_Lower & \texttt{effects(5) placebo(2) dont\_drop\_larger\_lower} \\
Effects\_Equal         & \texttt{effects(5) placebo(2) effects\_equal(all)} \\
Continuous\_1          & \texttt{effects(5) continuous(1)} \\
Continuous\_4          & \texttt{effects(5) continuous(4)} \\
\bottomrule
\end{tabular}
\end{table}

\subsection{Baseline}

Standard computation with no additional filtering. The working sample contains
all groups and periods after the standard data-cleaning pipeline. The
\texttt{avg\_time\_periods} block runs on the full cleaned dataset.

\begin{itemize}[nosep]
  \item \texttt{d1\_avg\_XX = d\_sq\_XX} (baseline binary treatment).
  \item \texttt{ngt\_bin\_XX = (outcome\_XX != .)}.
  \item $T_g^{(h)}$ computed as described in Section~5.
  \item $M_g = \min(5,\, T_g^{(h)} - F_g + 1)$ for each switcher.
\end{itemize}

\subsection{Placebos}

The \texttt{placebo()} option adds pre-switch ``placebo'' effect estimates at
lags $k = -1, -2, \dots$. The \texttt{avg\_time\_periods} computation
operates exclusively over \emph{post-switch} horizons ($k \geq 0$) and is
therefore completely unaffected by \texttt{placebo()}. All quantities are
identical to the Baseline case.

\subsection{Normalized}

\texttt{normalized} rescales each dynamic effect estimate by the average
magnitude of the treatment change $|D_{g,F_g} - D_{g,1}|$. The
\texttt{avg\_time\_periods} computation itself already weights by
$\Delta_{g,k} = |D_{g,\tau} - D_{g,1}|$ (the \texttt{incr} term in
Equation~\eqref{eq:avgM}). For binary union membership, $\Delta_{g,k} = 1$
for all post-switch observations of all switchers, so the weighting is
effectively uniform. The \texttt{normalized} flag has no effect on
\texttt{avg\_time\_periods}.

\subsection{Controls}

The \texttt{controls(hours)} option adds control variables to the DiD regression.
These enter via a pre-processing demeaning step that modifies the residualized
outcome used in the standard effect estimates (stored back into
\texttt{diff\_y\_i\_XX}).

The \texttt{avg\_time\_periods} computation uses \texttt{outcome\_XX} directly
(not the residualized outcome) for cell-validity checks
(\texttt{ngt\_bin\_XX = (outcome\_XX != .)}), and uses raw treatment-timing
variables ($F_g$, $T_g^{\text{cls}}$). Controls therefore have no effect on
$T_g^{(h)}$, $M_g$, the switcher counts, or the average $\widehat{\overline{M}}$.

\subsection{Trends\_Nonparam}

\texttt{trends\_nonparam(black)} stratifies the estimation by an additional
time-invariant group characteristic. It enters the data pipeline in two places
relevant to \texttt{avg\_time\_periods}:

\begin{enumerate}[nosep]
  \item \textbf{Class ceiling}: $T_g^{\text{cls}}$ is computed as
        $\max_{g' \in \text{cls}(g)} F_{g'}^{\text{trunc}} - 1$ where the class
        is now $\{g' : D_{g',1} = D_{g,1} \wedge \texttt{black}_{g'} = \texttt{black}_g\}$
        (line~1085). This may reduce $T_g^{\text{cls}}$ relative to the
        baseline, which in turn can reduce $T_g^{(h)}$ and $M_g$.

  \item \textbf{Control matching in the standard pipeline}: line~4744 computes
        control group counts within \texttt{time\_XX d\_sq\_XX trends\_nonparam}
        strata, so only same-strata controls contribute. In the Mata functions,
        the strata indicator passed as \texttt{CLS} is \texttt{d\_sq\_int\_XX}
        (which encodes $D_{g,1}$ only). On wagepan, where \texttt{black} is
        binary and creates no effective cross-strata confounds at the switcher
        counts level, the two definitions agree, which is why the spec passes.
\end{enumerate}

\subsection{Weight}

\texttt{weight(educ)} specifies a sampling/frequency weight. Weights enter the
standard pipeline via \texttt{N\_gt\_XX}, which includes the weight when
counting cells. In \texttt{\_avg\_cumul}, the per-effect switcher count is:

\[
  \texttt{nswitch}[k+1] \mathrel{+}= W_{r}
  \quad \text{where } W = \texttt{N\_gt\_XX}.
\]

So \texttt{e(N\_switch\_avg\_k)} is a \emph{weighted} count, consistent with
how \texttt{e(N\_switchers\_effect\_k)} is computed in the standard pipeline.
The numerator and denominator of $\widehat{\overline{M}}$ use the unweighted
treatment increment \texttt{incr = |dval - d1|}; weights affect only the
switcher counts.

\subsection{Cluster}

\texttt{cluster(hisp)} specifies a variable for cluster-robust standard errors.
Clustering affects only the variance estimation in Program~2; it does not enter
the point-estimate computation or the data-cleaning pipeline. In
\texttt{\_compute\_Tg} and \texttt{\_avg\_cumul}, the strata indicator
\texttt{CLS} is always \texttt{d\_sq\_int\_XX}, not the cluster variable.
Hence \texttt{avg\_time\_periods} is entirely unaffected by \texttt{cluster()},
and the switcher counts match the standard pipeline (which also does not use
the cluster variable in its control-count computation).

\subsection{Same\_Switchers}

\texttt{same\_switchers} restricts the estimator to use only groups that can
contribute to \emph{all} requested dynamic effects. In \texttt{\_avg\_cumul}
this is implemented via \texttt{ssw\_flag = 1}, which enforces the condition:

\[
  M_g \geq L \iff T_g^{(h)} - F_g + 1 \geq L \iff T_g^{(h)} \geq F_g + L - 1.
\]

Groups with fewer than $L$ valid post-switch periods are excluded from both the
average and the switcher counts. The standard pipeline enforces the analogous
restriction via the \texttt{still\_switcher\_i\_XX} indicator (line~4846/4870).

\subsection{Same\_Switchers\_Pl}

\texttt{same\_switch\-ers\_pl} further restricts the placebo estimates to the
same set of switchers used for the dynamic effects. The
\texttt{avg\_time\_periods} computation covers only post-switch effects, so
\texttt{same\_switch\-ers\_pl} has no additional effect beyond
\texttt{same\_switch\-ers}. The behaviour is identical to the
\textsc{Same\_Switchers} case.

\subsection{Switchers\_In}

\texttt{switchers(in)} restricts to groups whose average post-switch treatment
exceeds their baseline ($S_g = 1$, upward switchers). In \texttt{\_avg\_cumul},
the condition \texttt{sw\_dir != "in" or sg\_val == 1} excludes downward
switchers from the outer loop entirely. The standard pipeline enforces the same
restriction by running \texttt{did\_multiplegt\_dyn\_core\_new} with
\texttt{increase\_XX = 1} and checking \texttt{S\_g\_XX == increase\_XX} in
\texttt{distance\_to\_switch\_i\_XX}. Both sides restrict to the same set of
groups, so the counts agree.

\subsection{Switchers\_Out}

Symmetric to \textsc{Switchers\_In}: only groups with $S_g = 0$ (downward
switchers) contribute. \texttt{sw\_dir = "out"} in \texttt{\_avg\_cumul},
\texttt{increase\_XX = 0} in the standard pipeline.

\subsection{Only\_Never\_Switchers}

\texttt{only\_never\_switchers} restricts the control groups to true
never-switchers ($F_{g'} = T_{\max}+1$). In the standard pipeline this is
enforced at line~4738:

\begin{lstlisting}[language={}]
replace never_change_d_`i'_XX = 0
    if F_g_XX > time_XX & F_g_XX < T_max_XX + 1 & diff_y_`i'_XX != .
\end{lstlisting}

which zeroes out eventually-switching groups. In \texttt{\_compute\_Tg} and
\texttt{\_avg\_cumul}, the control check is $F_{g'} > P$ (i.e., hasn't
switched by period $P$). Never-switchers always satisfy $F_{g'} = T_{\max}+1
> P$; eventually-switching controls with a late $F_{g'}$ also pass this check
even under \texttt{only\_never\_switchers}. On wagepan, whenever an
eventually-switching group qualifies as a control in Mata, a true
never-switcher control is also available, so the switcher \emph{counts} agree
across both approaches (the test passes). The option does not affect
\texttt{avg\_time\_periods} inputs.

\subsection{CI\_Level\_90 and CI\_Level\_99}

\texttt{ci\_level()} controls the confidence level used when constructing
confidence intervals. It affects no part of the point-estimate or switcher-count
computation. Both variants are behaviorally identical to \textsc{Placebos} for
the purposes of \texttt{avg\_time\_periods}.

\subsection{Less\_Conservative\_SE}

\texttt{less\_conservative\_se} selects an alternative variance estimator.
Standard errors are estimated in Program~2 and do not feed back into the
\texttt{avg\_time\_periods} block. Completely unaffected.

\subsection{Dont\_Drop\_Larger\_Lower}

By default, groups that exhibit \emph{both} a strictly positive and a strictly
negative deviation from their baseline treatment at any point are dropped (line
886). \texttt{dont\_drop\_larger\_lower} suppresses this drop, retaining
``mixed-switcher'' groups in the sample. These groups are then assigned a
switch direction via $S_g = \mathbf{1}[\bar{D}_g^{\text{post}} > D_{g,1}]$
(where $\bar{D}_g^{\text{post}}$ is the average post-switch treatment).

For \texttt{avg\_time\_periods}:
\begin{itemize}[nosep]
  \item Mixed-switcher groups that are retained may now appear in the
        \texttt{\_avg\_cumul} loop.
  \item Their \texttt{incr}~$= |D_{g,\tau} - D_{g,1}|$ correctly uses the
        (possibly non-monotone) treatment path, consistent with the standard
        pipeline's use of \texttt{treatment\_XX}.
  \item $T_g^{(h)}$ for these groups is computed identically to other
        switchers.
\end{itemize}

On wagepan, \texttt{union} is binary so mixed-direction groups do not arise
(a group either stays at 0, moves to 1, or vice versa). The spec therefore
tests that the absence of such groups does not change behaviour.

\subsection{Effects\_Equal}

\texttt{effects\_equal(all)} adds a test of the joint hypothesis that all
dynamic effects are equal. The test is computed from the estimated coefficient
vector and its variance; it has no interaction with \texttt{avg\_time\_periods}.

%───────────────────────────────────────────────────────────────────────────────
\section{The \texttt{continuous()} Case}
%───────────────────────────────────────────────────────────────────────────────

\subsection{Continuous\_1 and Continuous\_4}

These two specs use a continuous treatment variable and request polynomial
controls of degree 1 and 4 respectively. They triggered the only
non-trivial implementation issue during development.

\subsubsection{How \texttt{continuous(p)} Changes the Dataset}

\begin{enumerate}[nosep]
  \item The original continuous baseline is saved: \texttt{d\_sq\_XX\_orig = d\_sq\_XX}.
  \item The baseline is replaced by zero: \texttt{d\_sq\_XX = 0}.
  \item Polynomial controls $D_{g,1}^j \times \sum_{k=1}^{t} \mathbf{1}\{t=k\}$
        for $j = 1,\dots,p$ are added to \texttt{controls}.
  \item The treatment is binarized:
        \[
          \tilde{D}_{g,t} =
          \begin{cases}
            \mathbf{1}[t \geq F_g] \cdot S_g^{\text{het}} & \text{if } S_g^{\text{het}} \neq \texttt{.}
                                                               \quad(\text{ever-switchers})\\
            \texttt{.} & \text{if } S_g^{\text{het}} = \texttt{.}
                          \quad(\text{never-switchers})
          \end{cases}
        \]
        where $S_g^{\text{het}} = +1$ for upward switchers and $-1$ for
        downward switchers.
\end{enumerate}

After step~4, \texttt{treatment\_XX = .} for every observation of a
never-switcher group.

\subsubsection{The Root Cause of the Offset Bug}

If the cell-validity indicator were defined as
\texttt{(outcome\_XX != .\ \&\ treatment\_XX != .)}, then every
never-switcher observation would have \texttt{ngt\_bin = 0}. When
\texttt{\_compute\_Tg} and \texttt{\_avg\_cumul} look for a control $g'$,
they require \texttt{NGT[$g'$, refp] == 1 and NGT[$g'$, $P$] == 1}. With
\texttt{ngt\_bin = 0} for never-switchers, these groups are invisible to Mata
as potential controls.

As a result, the only groups qualifying as controls are \emph{eventually-switching}
groups with $F_{g'} > P$, i.e., groups that switch later than the current
horizon. For wagepan, where the maximum $F_g$ among switchers equals $T_{\max}$:
\begin{itemize}[nosep]
  \item Mata counts a switcher $g$ as contributing to horizon $k$ only if
        there exists \emph{another switcher} with $F_{g'} > F_g + k$, giving
        $N_{\text{switch\_avg},k} = \#\{g : F_g \leq T_{\max} - k\}$.
  \item The standard pipeline counts a switcher $g$ at horizon $k$ if there
        exists \emph{any} control (including never-switchers) at $F_g + k - 1$,
        giving $N_{\text{std},k} = \#\{g : F_g \leq T_{\max} - k + 1\}$.
\end{itemize}

These two expressions satisfy $N_{\text{switch\_avg},k} = N_{\text{std},k+1}$:
a systematic one-step offset.

\subsubsection{The Fix}

Change the cell-validity indicator to depend only on the outcome:

\begin{lstlisting}[language={}]
-- gen byte `ngt_bin_XX' = (outcome_XX != . & treatment_XX != .)
++ gen byte `ngt_bin_XX' = (outcome_XX != .)
\end{lstlisting}

Never-switchers now have \texttt{ngt\_bin = 1} whenever their outcome is
observed (they always have valid outcomes on the balanced wagepan panel). The
Mata functions find them as valid controls, restoring
$N_{\text{switch\_avg},k} = N_{\text{std},k}$ for all $k$.

\subsubsection{Baseline Treatment for the Average Formula}

For \texttt{continuous(p)}, the treatment magnitude is computed using the
\emph{original} continuous baseline:

\begin{lstlisting}[language={}]
local d1_avg_XX "d_sq_XX_orig"    // not d_sq_XX, which is 0 after line 923
\end{lstlisting}

So $\Delta_{g,k} = |D_{g,\tau}^{\text{orig}} - D_{g,1}^{\text{orig}}|$ uses
the pre-binarization continuous values. This is consistent with the standard
pipeline's use of \texttt{treatment\_XX\_orig} when computing treatment
magnitudes under \texttt{continuous()}.

\subsubsection{Polynomial Degree}

The degree $p$ of the polynomial controls (\texttt{continuous(1)} vs
\texttt{continuous(4)}) determines how many interaction controls are added to
the regression. It affects the \emph{estimator} values but not the
cell-validity checks, timing variables, or control-group eligibility used by
\texttt{avg\_time\_periods}. Both degrees therefore produce the same switcher
counts and the same $\widehat{\overline{M}}$ structure, and both pass the
cross-check.

%───────────────────────────────────────────────────────────────────────────────
\section{Summary of Option Interactions}
%───────────────────────────────────────────────────────────────────────────────

\begin{table}[ht]
\centering
\caption{How each option interacts with \texttt{avg\_time\_periods}}
\label{tab:interactions}
\renewcommand{\arraystretch}{1.25}
\begin{tabular}{@{}lp{2.8cm}p{8cm}@{}}
\toprule
Option & Affects \texttt{avg\_tp}? & Mechanism \\
\midrule
\texttt{placebo()} & No & Pre-switch only; \texttt{avg\_tp} is post-switch \\
\texttt{normalized} & No & Rescales estimates, not timing or counts \\
\texttt{controls()} & No & Residualizes outcomes; \texttt{avg\_tp} uses raw \texttt{outcome\_XX} \\
\texttt{trends\_nonparam()} & Yes & Modifies $T_g^{\text{cls}}$ (class ceiling) via within-strata max of $F_g^{\text{trunc}}$ \\
\texttt{weight()} & Yes (counts) & \texttt{N\_switch\_avg\_k} is weighted by \texttt{N\_gt\_XX} \\
\texttt{cluster()} & No & SE only; strata for Mata is always \texttt{d\_sq\_int\_XX} \\
\texttt{same\_switchers} & Yes & \texttt{ssw\_flag=1}: excludes $g$ with $M_g < L$ \\
\texttt{same\_switchers\_pl} & No & Placebo restriction; \texttt{avg\_tp} unaffected \\
\texttt{switchers(in/out)} & Yes & Restricts to upward/downward switchers in outer loop \\
\texttt{only\_never\_switchers} & No$^*$ & Control eligibility differs but counts agree on wagepan \\
\texttt{ci\_level()} & No & CI only \\
\texttt{less\_conservative\_se} & No & SE only \\
\texttt{dont\_drop\_larger\_lower} & Yes (sample) & Retains mixed-direction groups in sample \\
\texttt{effects\_equal()} & No & Aggregation test only \\
\texttt{continuous(p)} & Yes & Binarizes treatment; uses \texttt{d\_sq\_XX\_orig} for $\Delta$; requires \texttt{ngt\_bin = (outcome != .)} \\
\bottomrule
\end{tabular}
\vspace{4pt}

{\footnotesize $^*$ On wagepan, eventually-switching controls are always
accompanied by a never-switcher control, so the restriction is vacuous for
the switcher \emph{counts}. The control check in Mata ($F_{g'} > P$) does not
explicitly enforce the \texttt{only\_never\_switchers} restriction.}
\end{table}

%───────────────────────────────────────────────────────────────────────────────
\section{Verification Procedure}
%───────────────────────────────────────────────────────────────────────────────

The invariant $\texttt{e(N\_switch\_avg\_k)} = \texttt{e(N\_switchers\_effect\_k)}$
for all $k$ and all wagepan specifications is verified by:

\begin{enumerate}[nosep]
  \item \textbf{Stata}: \texttt{stata-se -b do "code/test\_wagepan\_specs.do"}\\
        Runs all 19 specs, captures both counts per effect, writes
        \texttt{logfiles/wagepan\_specs.csv}.
  \item \textbf{Python}: \texttt{python3 code/compare\_wagepan\_specs.py}\\
        Reads the CSV, reports PASS/FAIL/MISS per spec and per effect,
        exits with code~1 if any discrepancy is found.
\end{enumerate}

Expected output: all 19 specs with PASS on all 5 effects each (95 rows total).

\end{document}
